\documentclass[11pt]{amsart}
\usepackage{amssymb,amsmath}
\usepackage[margin=3cm]{geometry}
\usepackage{booktabs}

\newcommand{\stamptime}{\the\numexpr(\time-30)/60\relax:\ifnum\numexpr\time-60*((\time-30)/60)\relax<10 0\fi\the\numexpr\time-60*((\time-30)/60)\relax}
\AddToHook{shipout/foreground}{\put(0.5\paperwidth,-\paperheight+0.55cm){\makebox[0pt]{\tiny\ttfamily unit/doc/\jobname.tex \ --- \ compiled \today, \stamptime}}}

\newtheorem{theorem}{Theorem}[section]
\newtheorem{lemma}[theorem]{Lemma}
\newtheorem{proposition}[theorem]{Proposition}
\newtheorem{corollary}[theorem]{Corollary}
\theoremstyle{definition}
\newtheorem{definition}[theorem]{Definition}
\theoremstyle{remark}
\newtheorem{remark}[theorem]{Remark}

\newcommand{\seq}[1]{\langle #1\rangle}
\newcommand{\one}[1]{[#1]}
\newcommand{\nx}{\bar x}
\newcommand{\N}{\mathbb{N}}

\title{Novikov's consistency proof}
\author{}
\date{}

\begin{document}
\begin{abstract}
  In 1941 Novikov proved, intuitionistically, that the rule of inference
  is admissible for an infinitary propositional calculus~\cite{Novikov}. As
  consequences he obtained the consistency of arithmetic and a method for
  extracting a witness from a classical proof of an existential statement.
  We present his argument in the form of the appendix
  of~\cite{Coquand}. The cut rule is proved by induction on regularity
  proofs. As Mints stressed~\cite{Mints}, there is no induction on the
  cut formula. The proofs below
  have been checked in the proof checker Unit (Section~\ref{sec:unit}).
\end{abstract}
\maketitle

\section{Formulas and sequents}

We fix once and for all a truth value for each of the constants $x_n$,
$n\in\N$: either $x_n$ is true and $\nx_n$ is false, or the converse. This
is Novikov's choice of the $x'_n$~\cite[\S3]{Novikov}.

\begin{definition}
  The \emph{formulas} are inductively generated by $0$, $1$, $x_n$,
  $\nx_n$, and $\seq{A_n}$ for every sequence $(A_n)_{n\in\N}$ of formulas.
\end{definition}

A formula has two readings. Read as a \emph{sum}, $\seq{A_n}$ is
$\sum_n A_n$ with each $A_n$ read as a product. Read as a \emph{product},
$\seq{A_n}$ is $\prod_n A_n$ with each $A_n$ read as a sum. The formulas $0$,
$1$, $x_n$, $\nx_n$ have the same meaning in both readings, and $0$ and $1$
are the empty sum and the empty product. This is Novikov's normal form
(\S6): the summands of a sum are products, and the factors of a product
are sums. Finite sums and products are infinite ones with repeated terms,
and Novikov's $\vee$, $\cdot$, $\to$ reduce to these.

The $k$-th \emph{summand} of a formula read as a sum is $A_{(k)}$, where
$\seq{A_n}_{(k)} = A_k$ and $c_{(k)}=c$ for $c$ one of $0,1,x_n,\nx_n$. The
$k$-th \emph{factor} of a formula read as a product is the same
$A_{(k)}$. \emph{Negation} is the de Morgan dual: $\neg 0=1$, $\neg 1=0$,
$\neg x_n=\nx_n$, $\neg\nx_n=x_n$, $\neg\seq{A_n}=\seq{\neg A_n}$. In
particular $(\neg A)_{(k)}=\neg(A_{(k)})$. For a formula $P$ read as a
product, $\one{P}=\seq{P,P,\dots}$ is the sum whose only summand is $P$.

\begin{definition}
  A \emph{sequent} $\Gamma=\Gamma_0,\Gamma_1,\dots$ is a sequence of
  formulas read as sums, a finite sequent being completed by $0$'s. It
  stands for the sum of all the summands $\Gamma_{m(k)}$, which is
  Novikov's merged sum $\widetilde{\sum}$ (\S3). We write $A,\Gamma$ for
  the sequent with $A$ added in front, and $\emptyset$ for the empty
  sequent. The \emph{negation} of a sequent is the product
  $\neg\Gamma=\seq{\one{\neg\Gamma_m}}_m$. Its factors are the sums
  $\one{\neg\Gamma_m}$, whose only summand is $\neg\Gamma_m$, read as a
  product.
\end{definition}

A \emph{false atom} is $0$ or a false constant.

\section{Regular sequents}

\begin{definition}[Novikov, \S4]
  The regular sequents are defined inductively. $\Gamma$ is regular if
  it has a summand $c$ such that
  \begin{enumerate}
  \item $c$ is $1$ or a true constant, or
  \item $c=\seq{C_j}$ and $C_j,\Gamma$ is regular for every $j$.
  \end{enumerate}
  The summand $c$ is the \emph{distinguished term}, and the sequents
  $C_j,\Gamma$ are the \emph{constituents}.
\end{definition}

The definition is an inductive one, so we can prove a property of all
regular sequents by showing that it holds of $\Gamma$ whenever it holds of
all the constituents. This is \emph{induction on regularity}. Novikov
removes the distinguished term from the constituents. Here, as in Tait's
calculus, it is kept. By Lemma~\ref{lem:incl}, the two definitions give the
same regular sequents.

Two observations follow at once.

\begin{lemma}\label{lem:empty}
  The empty sequent is not regular. More generally, a sequent all of whose
  summands are false atoms is not regular.
\end{lemma}

\begin{proof}
  No summand can be the distinguished term.
\end{proof}

\begin{lemma}[Inclusion; Novikov's Lemmas 1, 3, 7, 14]\label{lem:incl}
  Suppose every summand of $\Gamma$ is a false atom or a summand of $\Delta$. If
  $\Gamma$ is regular, so is $\Delta$.
\end{lemma}

\begin{proof}
  By induction on regularity of $\Gamma$, for all $\Delta$. The
  distinguished term $c$ is not a false atom, so it is a summand of $\Delta$. If
  $c=\seq{C_j}$, the inclusion hypothesis also holds for $C_j,\Gamma$ and
  $C_j,\Delta$, so the induction hypothesis gives that $C_j,\Delta$ is
  regular. Hence $\Delta$ is regular, with the same distinguished term.
\end{proof}

The lemma covers weakening, exchange and contraction. It also covers the
removal of summands that are false atoms, and replacing $\one{P}$ by a sum that has $P$
among its summands.

\begin{lemma}[Identity; Novikov \S8]\label{lem:id}
  For every formula $A$ and sequent $\Gamma$, the sequents
  $A,\one{\neg A},\Gamma$ and $\one{A},\neg A,\Gamma$ are regular.
\end{lemma}

\begin{proof}
  By induction on $A$, the two statements together. For a constant, the
  true one of $x_n,\nx_n$ is the distinguished term, and for $0$ and $1$
  the summand $1$ is. Let $A=\seq{A_j}$.

  For the first statement, take the distinguished term $\neg A$, the
  summand of $\one{\neg A}$. Its constituents are
  $\neg A_j,A,\one{\neg A},\Gamma$. By the induction hypothesis
  $\one{A_j},\neg A_j,\Gamma$ is regular, and $A_j$ is a summand of $A$.
  So the constituents are regular by Lemma~\ref{lem:incl}.

  For the second statement, take the distinguished term $A$, the summand
  of $\one{A}$. Its constituents $A_i,\one{A},\neg A,\Gamma$ include
  $A_i,\one{\neg A_i},\Gamma$, since $\neg A_i$ is a summand of the sum
  $\neg A$. That sequent is regular by the induction hypothesis.
\end{proof}

This is the only induction on formulas in the proof of the cut rule. It
is applied only to the formulas of the context (Theorem~\ref{thm:cut}).

The remaining lemmas have a common form. If $X,B$ is regular, then so is
$Y,B$, for all $B$. Such a lemma is proved by induction on a proof of a
sequent $\Gamma$ whose summands are false atoms or summands of $X,B$, for all $B$.
If the distinguished term is a summand of $B$, the same step applies to
$Y,B$. Indeed, by the induction hypothesis (with $C_j,B$ for $B$) the
sequent $Y,C_j,B$ is regular, and its summands are those of $C_j,Y,B$.
Only the case of a distinguished term among the summands of $X$ needs an
argument. We call this \emph{induction on $X,B$}.

\begin{lemma}[Inversion; Novikov's Lemma 1]\label{lem:inv}
  If $\one{P},B$ is regular, then so is $P_{(i)},B$ for every $i$.
\end{lemma}

\begin{proof}
  By induction on $\one{P},B$, where the distinguished term is $P$. If $P$
  is $1$ or a true constant, then $P_{(i)}=P$ makes $P_{(i)},B$ regular. If
  $P=\seq{C_j}$, then the constituent $C_i,\Gamma$ is regular, and the
  induction hypothesis (with $C_i,B$ for $B$) gives that $C_i,C_i,B$ is
  regular. Lemma~\ref{lem:incl} then gives $C_i,B$.
\end{proof}

\section{The rule of inference}

Following the appendix of~\cite{Coquand}, we call a sequent $D$
\emph{eliminable} if, for every sequent $B$, $B$ is regular whenever
$\one{\neg D},B$ is regular.

\begin{lemma}[Lemma 5 of~\cite{Coquand}; Novikov's Lemmas 9--12]\label{lem:and}
  Let $D$ be a sequent and $(e_j)$ a sequence of formulas. If
  $\one{\neg D},B$ and $\seq{\neg e_j}_j,B$ are regular, then so is
  $Z,B$ with $Z=\seq{\neg(e_j,D)}_j$.
\end{lemma}

The summand $\neg(e_j,D)$ of $Z$ is the product of $\one{\neg e_j}$ and of
the $\one{\neg D_m}$. In other words, the product $\neg D$ is moved inside
each summand $\neg e_j$ of the sum.

\begin{proof}
  By induction on $X,B$ with $X=\seq{\neg e_j}_j$, for all $B$ such that
  $\one{\neg D},B$ is regular. This property of $B$ is preserved when
  formulas are added to $B$, by Lemma~\ref{lem:incl}. Let the
  distinguished term be $\neg e_l$. We show that $Z,B$ is regular, with
  distinguished term $\neg(e_l,D)$. Its constituents are of two kinds.
  \begin{itemize}
  \item $\one{\neg D_m},Z,B$. By inversion (Lemma~\ref{lem:inv}) of
    $\one{\neg D},B$, the sequent $\one{\neg D_m},B$ is regular, and so is
    this constituent by Lemma~\ref{lem:incl}.
  \item $\one{\neg e_l},Z,B$. We take $\neg e_l$ as distinguished term. If
    it is $1$ or a true constant, we are done. If $\neg e_l=\seq{C_i}$,
    the constituents $C_i,\Gamma$ of the given proof satisfy the induction
    hypothesis, which gives that $Z,C_i,B$ is regular. This sequent is
    included in the constituent $C_i,\one{\neg e_l},Z,B$ needed here.\qedhere
  \end{itemize}
\end{proof}

\begin{lemma}[Lemma 6 of~\cite{Coquand}]\label{lem:sum}
  If every $D_j$ is eliminable and $\seq{\neg D_j}_j,B$ is regular, then
  $B$ is regular.
\end{lemma}

\begin{proof}
  By induction on $X,B$ with $X=\seq{\neg D_j}_j$ and $Y=0$. Let the
  distinguished term be $\neg D_l=\seq{\one{\neg D_{l,m}}}_m$. Its
  constituents $\one{\neg D_{l,m}},\Gamma$ give, by the induction
  hypothesis, that $\one{\neg D_{l,m}},B$ is regular for every $m$. These
  sequents give the constituents of $\one{\neg D_l},B$ with
  distinguished term $\neg D_l$, so $\one{\neg D_l},B$ is regular. Since
  $D_l$ is eliminable, $B$ is regular.
\end{proof}

\begin{proposition}[Proposition 2 of~\cite{Coquand}]\label{prop:elim}
  Every regular sequent is eliminable.
\end{proposition}

\begin{proof}
  By induction on regularity of $D$. Let $c=D_{r(k)}$ be the distinguished
  term of $D$, and let $\one{\neg D},B$ be regular. We prove that $B$ is
  regular by induction on $\one{\neg D},B$ (with $Y=0$). The distinguished
  term is $\neg D$, with constituents $\one{\neg D_m},\Gamma$. The
  induction hypothesis gives that $\one{\neg D_m},B$ is regular for every
  $m$, in particular for $m=r$. Inversion at $k$ then gives that
  $\neg c,B$ is regular, since $(\neg D_r)_{(k)}=\neg c$.
  \begin{itemize}
  \item If $c$ is $1$ or a true constant, then $\neg c$ is a false atom, and $B$
    is regular by Lemma~\ref{lem:incl}.
  \item If $c=\seq{g_j}$, then $\neg c=\seq{\neg g_j}$, and Lemma~\ref{lem:and}
    with $e=g$ gives that $\seq{\neg(g_j,D)}_j,B$ is regular. The sequents
    $g_j,D$ are the constituents of $D$, so they are eliminable by the
    induction hypothesis on $D$. Lemma~\ref{lem:sum} gives that $B$ is
    regular.\qedhere
  \end{itemize}
\end{proof}

\begin{remark}
  No induction on the cut formula is used. The outer induction is on the
  regularity proof of $D$, and its induction hypothesis is applied to the
  constituents $g_j,D$, which as formulas are larger than $D$. The inner
  inductions are on the regularity proof of the other premise. The shape
  of the distinguished term $c$ only gives a case distinction. Mints
  insisted on this feature of Novikov's argument~\cite{Mints},
  which distinguishes it from Gentzen's proof. Tupailo compared the two
  methods of cut elimination~\cite{Tupailo}.
\end{remark}

\begin{theorem}[Cut]\label{thm:cut}
  If $A,\Gamma$ and $\one{\neg A},\Gamma$ are regular, then so is
  $\Gamma$.
\end{theorem}

\begin{proof}
  By Proposition~\ref{prop:elim}, the sequent $D=A,\Gamma$ is eliminable.
  It remains to show that $\one{\neg D},\Gamma$ is regular, with
  distinguished term $\neg D$. Its constituents are
  $\one{\neg A},\one{\neg D},\Gamma$, which is regular by hypothesis and
  Lemma~\ref{lem:incl}, and $\one{\neg\Gamma_m},\one{\neg D},\Gamma$,
  which includes the regular sequent $\Gamma_m,\one{\neg\Gamma_m}$ of
  Lemma~\ref{lem:id}.
\end{proof}

\begin{corollary}[Novikov's rule of inference]
  If $A$ and $\one{\neg A},\Gamma$ are regular, then so is $\Gamma$. The
  formulas $A$ and $\one{\neg A}$ are not both regular.
\end{corollary}

\begin{proof}
  By Lemma~\ref{lem:incl}, $A,\Gamma$ is regular, and Theorem~\ref{thm:cut}
  applies. For the second statement take $\Gamma=\emptyset$ and use
  Lemma~\ref{lem:empty}.
\end{proof}

Novikov's ``true formulas'' (\S3) are the theorems of a Hilbert-style
calculus. Its axioms $\prod_iF_i\to F_n$, $F_n\to\sum_iF_i$ and its rule
6 are instances of the definition of regularity, and its rule of inference
is the Corollary. Equivalently, consider the sequent calculus whose rules
are the two clauses of the definition of regularity together with the cut
rule: from $A,\Gamma$ and $\one{\neg A},\Gamma$ infer $\Gamma$. By
induction on derivations and Theorem~\ref{thm:cut}, \emph{every derivable
sequent is regular}, which is Novikov's principal theorem (\S8). By
Lemma~\ref{lem:empty}, \emph{the calculus is consistent}.

\subsection*{Comparison with Novikov's proof}

Novikov's proof of the rule of inference (\S\S5--8) has the same two
inductions, organised globally.

\medskip
\begin{tabular}{@{}p{0.55\textwidth}p{0.38\textwidth}@{}}
\toprule
Novikov~\cite{Novikov} & here, after~\cite{Coquand} \\
\midrule
The symbol $\tau$ is built from the regularity proof of $F$. The operator
$B_\tau$ turns $F$ into the product of the elementary regular formulas at
the leaves of that proof (\S7). &
the induction on regularity of $D$ (Proposition~\ref{prop:elim}) \\
\addlinespace
Lemmas 9--12, the opening of brackets, used by $A_\tau$ in Lemmas 17, 18. &
Lemma~\ref{lem:and} \\
\addlinespace
Lemma 15, by induction on regularity of $\sum_iA_{\tau_i}(\bar F_i)\vee G$,
removes the pairs $A\bar A$ (atomic cuts, Lemma 13) and the false
constants (Lemma 14). &
Lemma~\ref{lem:sum} and the inner induction of
Proposition~\ref{prop:elim} \\
\bottomrule
\end{tabular}
\medskip

The two proofs have the same inductions and the same ingredients,
organised differently. Novikov first unfolds the whole regularity proof of
$F$ into its axioms and transforms $\bar F\vee G$ in one step. He then
removes the duals of the axioms: the pairs $A\bar A$ on letters by atomic
cuts (Lemma 13) and the false constants (Lemma 14). The proof above
proceeds by induction on the regularity proof of $D$, using
``eliminable'' as the induction hypothesis. Since the formulas here are
closed, only the analogue of Lemma 14 is needed, namely the removal of
false atoms. We do not claim that the two proofs compute the same
cut-free proof.

\section{Arithmetic}

\subsection*{Induction (Novikov \S9)}
Let $A_0,A_1,\dots$ be formulas read as sums. The induction principle
$A_0\wedge\forall i\,(A_i\to A_{i+1})\to\forall i\,A_i$ is the sequent
\[
  S \;=\; \one{\neg A_0},\ \seq{A_i\wedge\one{\neg A_{i+1}}}_i,\ \one{\seq{A_i}},
\]
where $A\wedge B$ is the product $\seq{A,B,B,\dots}$.

\begin{proposition}
  $S$ is regular, and its proof uses no cut.
\end{proposition}

\begin{proof}
  Take the distinguished term $\seq{A_i}$. We show that each constituent
  $A_n,S$ is regular, by induction on $n$. For $n=0$ this is
  Lemma~\ref{lem:id}. For $n+1$, take the distinguished term
  $A_n\wedge\one{\neg A_{n+1}}$. Its constituents are
  $A_n,A_{n+1},S$, which follows from the case $n$ by Lemma~\ref{lem:incl},
  and $\one{\neg A_{n+1}},A_{n+1},S$, which follows from
  Lemma~\ref{lem:id}.
\end{proof}

Arithmetic with primitive recursive functions is interpreted with
$\forall=\prod$ and $\exists=\sum$, a closed atomic formula $t=s$ being
the constant whose truth is that of the equation. Under this
interpretation the axioms are regular, and the rules of inference
preserve regularity. Hence arithmetic is consistent~\cite[\S9]{Novikov}.

\subsection*{Existence proofs (Novikov \S10)}

A \emph{finitary formula} is built from the constants with $\wedge$ and
$\vee$. Its truth is computed from the truth of the constants. We read it
as a product $\ulcorner F\urcorner$: a constant is itself,
$\ulcorner F\wedge G\urcorner=\seq{\one{\ulcorner F\urcorner},
\one{\ulcorner G\urcorner},\dots}$, and
$\ulcorner F\vee G\urcorner=\one{\seq{\ulcorner F\urcorner,\ulcorner
G\urcorner,\dots}}$.

\begin{lemma}\label{lem:dec}
  If $F$ is true, then $\one{\ulcorner F\urcorner},\Gamma$ is regular. If
  $F$ is false, then $\neg\ulcorner F\urcorner,\Gamma$ is regular.
\end{lemma}

\begin{proof}
  By induction on $F$, the two statements together, taking the
  distinguished term at the top. For $F\vee G$ with $F$ true, the sequent
  $\one{\ulcorner F\urcorner},\Gamma$ is included in the constituent.
\end{proof}

\begin{theorem}[Novikov \S10; Markov's rule]
  Let $F(0),F(1),\dots$ be finitary formulas. If $\sum_n\ulcorner
  F(n)\urcorner$ is regular, in particular if it is derivable in the
  calculus with cut, then one can find $n$ such that $F(n)$ is true.
\end{theorem}

\begin{proof}
  We follow Novikov. Let $x$ be a true constant, and let $y_n=x$ if $F(n)$
  is true and $y_n=\neg x$ otherwise. The implication
  $\one{\neg\sum_n\ulcorner F(n)\urcorner},\seq{y_n}$ is regular, with
  distinguished term $\prod_n\neg\ulcorner F(n)\urcorner$. Each
  constituent $\neg\ulcorner F(n)\urcorner,\dots$ is regular: if $F(n)$
  is true, take $y_n$ as distinguished term; if $F(n)$ is false, use
  Lemma~\ref{lem:dec}. By the Corollary, $\seq{y_n}$ is regular. A
  summand of $\seq{y_n}$ is a constant, so the distinguished term of a
  proof of $\seq{y_n}$ is a true $y_n$. Then $y_n=x$, and $F(n)$ is true.
\end{proof}

The number $n$ is read off the regularity proof of $\seq{y_n}$, which the
proof of Theorem~\ref{thm:cut} computes from the given proof of
$\sum_n\ulcorner F(n)\urcorner$.

\section{The formalisation in Unit}\label{sec:unit}

The proofs above are checked in Unit (\texttt{github.com/coquand/unit}),
directory \texttt{novikov/}: 212 declarations in 0.02\,s. There are no
axioms. The truth values of the constants are the one parameter,
$\mathit{val}:\N\to\mathrm{Bool}$.
\begin{itemize}
\item \texttt{Formula}: formulas
  \texttt{F0 | F1 | X n | NX n | Sp f}, negation, summands, sequents as
  sequences $\N\to T$, and finitary formulas.
\item \texttt{Regular}: regularity is an inductive family
  \texttt{Reg} indexed by sequents, with one constructor for each clause
  of the definition. Induction on regularity is its eliminator. The file
  also contains Lemma~\ref{lem:empty}, Lemma~\ref{lem:incl},
  and the scheme of induction on $X,B$ as one module (\texttt{Replace}).
\item \texttt{Identity}: Lemmas~\ref{lem:id} and~\ref{lem:inv}.
\item \texttt{Cut}: Lemmas~\ref{lem:and} and~\ref{lem:sum},
  Proposition~\ref{prop:elim}, Theorem~\ref{thm:cut} and the Corollary.
\item \texttt{Calculus}: the calculus with cut, as an inductive family
  with a third constructor for the cut rule; every derivable sequent is
  regular, and the calculus is consistent.
\item \texttt{Arithmetic}, \texttt{Existence}: induction, and Novikov's
  existence theorem.
\item \texttt{Examples}: with $x_n$ true exactly for even $n$ and
  $F(n)=\nx_n\wedge x_{n+1}$, a proof of $\sum_nF(n)$ by its summand $F(3)$
  is given. The checker evaluates the witness through the elimination of
  the cut and obtains~$3$.
\end{itemize}

\begin{thebibliography}{9}
\bibitem{Novikov} P. S. Novikoff, \emph{On the consistency of certain
  logical calculus}, Recueil Math\'ematique (Mat.\ Sbornik) 12 (54)
  (1943), 231--261.
\bibitem{Coquand} T. Coquand, \emph{A semantics of evidence for classical
  arithmetic}, Journal of Symbolic Logic 60 (1995), 325--337. Appendix:
  Novikoff's proof.
\bibitem{Mints} G. Mints, \emph{Proof theory in the USSR 1925--1969},
  Journal of Symbolic Logic 56 (1991), 385--424.
\bibitem{Tupailo} S. Tupailo, \emph{Gentzen-style and Novikov-style
  cut-elimination}, in: A. Nerode, M. Taitslin (eds.), Logical Foundations
  of Computer Science (Tver '92), Lecture Notes in Computer Science 620,
  Springer, 1992, 493--502.
\bibitem{Beklemishev} L. D. Beklemishev, \emph{Novikov's work in proof
  theory: method of regular formulas}, lecture, 27 August 2026.
\bibitem{Tait} W. W. Tait, \emph{Normal derivability in classical logic},
  in: The Syntax and Semantics of Infinitary Languages, Lecture Notes in
  Mathematics 72, Springer, 1968, 204--236.
\end{thebibliography}

\end{document}
